\documentclass[12pt,a4paper,sans,fontset=windows]{moderncv} % Font sizes: 10, 11, or 12; paper sizes: a4paper, letterpaper, a5paper, legalpaper, executivepaper or landscape; font families: sans or roman
%\usepackage{xeCJK}
%\usepackage[utf8]{inputenc}
%\usepackage[export]{adjustbox}
\usepackage{amssymb,amsmath,amsthm}
\usepackage{times} %字体设置
\usepackage{fontspec}
%\setmainfont{Times New Roman}
\usepackage{graphicx}
\usepackage{wrapfig}
\usepackage{multicol}
\usepackage{enumitem}
%\usepackage{ctex}

\usepackage{xeCJK}
%\setCJKmainfont{SimSun}  % Windows 上的宋体
\setCJKmainfont{STSong}  % macOS 上的宋体
%中文需要用xelatex编译

\moderncvstyle{casual} % CV theme - options include: 'casual' (default), 'classic', 'oldstyle' and 'banking'
\moderncvcolor{blue} % CV color - options include: 'blue' (default), 'orange', 'green', 'red', 'purple', 'grey' and 'black'
\usepackage{color}
%\usepackage{lipsum} % Used for inserting dummy 'Lorem ipsum' text into the template

\usepackage[scale=0.8]{geometry} % Reduce document margins
\setlength{\hintscolumnwidth}{90pt} % Uncomment to change the width of the dates column 3cm 95pt
%\setlength{\makecvtitlenamewidth}{10cm} % For the 'classic' style, uncomment to adjust the width of the space allocated to your name

%----------------------------------------------------------------------------------------
%	NAME AND CONTACT INFORMATION SECTION
%----------------------------------------------------------------------------------------

\firstname{Yu} % Your first name
\familyname{Zhang} % Your last name

% All information in this block is optional, comment out any lines you don't need
\title{CV}
%\address{14, Gonghangdae-ro 57gil, Gangseo-gu, Seoul, Republic of Korea}
%\mobile{+82 10 8635 3011}
%\phone{(000) 111 1112}
%\fax{(000) 111 1113}
%\email{jayeon0724@gmail.com}
%\homepage{staff.org.edu/~jsmith}{staff.org.edu/$\sim$jsmith} % The first argument is the url for the clickable link, the second argument is the url displayed in the template - this allows special characters to be displayed such as the tilde in this example
%\extrainfo{additional information}
%\photo[70pt][0.4pt]{pictures/picture} % The first bracket is the picture height, the second is the thickness of the frame around the picture (0pt for no frame)
%\quote{}

%----------------------------------------------------------------------------------------

\begin{document}

\begin{center}
\vspace{10pt}
\textit{\Huge{\textcolor{black}{张宇}}}
\end{center}

\vspace{10pt} 

\section{工作经历}

\vspace{3pt} 

\cvitem{2023.09 -- ~ 至今~ }{\textbf{讲师}~ ~  \quad 天津大学 \quad  \textit{天津, 中国} }

\vspace{3pt} 

\cvitem{2020.09 -- 2023.09}{\textbf{博士后} \quad  南开大学 \quad  \textit{天津, 中国} }

\vspace{5pt} 


\section{教育经历}

\vspace{3pt} 

\cvitem{2014.08 -- 2020.05}{俄亥俄州立大学(The Ohio State University) \quad \textit{哥伦布, 俄亥俄, 美国} 
\newline 数学博士,  导师: John E. Harper}

\vspace{3pt} 

\cvitem{2010.09 -- 2014.07}{\textbf{北京大学} \quad \textit{北京, 中国}
\newline 数学学士, 导师: 范辉军}

\vspace{5pt} 


\section{科研基金}

\vspace{3pt} 

\cvitem{2025.01 -- 2027.12}{国家自然科学基金青年项目(No. 12401085): 基于motivic理论与计算机辅助计算的球面稳定同伦群奇质数分量群研究; \textbf{项目负责人}。}

\vspace{3pt} 

\cvitem{2024.01 -- 2025.12}{天津大学科技创新领军人才培育计划启明项目(No. 2024XQM-0009)：球面稳定同伦群的计算；\textbf{项目负责人}。}


\vspace{3pt} 

\cvitem{2023.01 -- 2026.12}{国家自然科学基金地区科学基金项目(No. 12261091)：有关谱序列及分类空间$BPU_n$的上同调群的若干问题研究；\textbf{合作单位项目负责人}。}

\vspace{3pt} 

\cvitem{2023.01 -- 2026.12}{国家自然科学基金面上项目(No. 12271183)：环面空间的上同调与motivic稳定同伦；\textbf{主要参与者}。}

\vspace{3pt} 

\cvitem{2021.05 -- 2023.04}{博士后国际交流计划引进项目： 幂零环结构谱的性质研究；\textbf{项目负责人}。}


\section{学术论文}

\vspace{3pt} 

\begin{itemize}[wide=30pt, leftmargin=*]
    \item \textbf{The secondary periodic element $\beta_{p^2/p^2-1}$ and its applications}
        \newline (J. Hong, X. Wang and \textbf{Y. Zhang}), \textit{SCIENCE CHINA Mathematics}. 68, 207-222 (2025).
    \item \textbf{Detecting nontrivial products in the stable homotopy ring of spheres via the third Morava stabilizer algebra}
    \newline (X. Wang, J. Wu, \textbf{Y. Zhang} and L. Zhong), \textit{Proceedings of the American Mathematical Society}. 152 (2024), 4521-4536.
    \item \textbf{A correspondence between higher Adams differentials and higher algebraic Novikov differentials at odd primes}
        \newline (X. Wang and \textbf{Y. Zhang}), \textit{Proceedings of the American Mathematical Society}. 151 (2023), 5087-5096.
    \item \textbf{The $p$-primary subgroup of the cohomology of $BPU_n$ in dimension $2p+6$}
        \newline (\textbf{Y. Zhang}, Z. Zhang and L. Zhong), \textit{Topology and Its Applications}.  338 (2023): 108642.
    \item \textbf{Some nontrivial secondary Adams differentials on the fourth line}
        \newline (X. Wang, Y. Wang and \textbf{Y. Zhang}),  \textit{New York Journal of Mathematics}.  29 (2023) 687-707.
    \item \textbf{Homotopy pro-nilpotent structured ring spectra and topological Quillen localization}
        \newline (\textbf{Y. Zhang}), \textit{Journal of Homotopy and Related Structures}. 17(4), 511-523 (2022).
    \item \textbf{The $p$-primary subgroups of the cohomology of $BPU_n$ in dimensions less than $2p+5$}
        \newline (X. Gu, \textbf{Y. Zhang}, Z. Zhang, and L. Zhong), \textit{Proceedings of the American Mathematical Society}. 150(9), 4099-4111 (2022).         
\end{itemize}

\vspace{2pt} 


\section{教育教学}

\vspace{2pt}

\cvitem{25261学期至今}{担任8位本科生的师友导师}

\vspace{2pt}

\cvitem{25261-2学期}{指导一名本科生毕业论文}

\vspace{2pt}

\cvitem{24251学期}{微积分I \quad 学生162人, 共96学时。
\newline 学生给出的量化评教平均分为: 97.6 (满分100)。}

\vspace{10pt}
 
\cvitem{}{\textit{在南开大学作为主讲老师授课：}}

\vspace{2pt}

\cvitem{21222学期}{文科概率统计 \quad 学生141人, 共48学时。
\newline 学生给出的量化评教平均分为: 97.7 (满分100)。}

\vspace{2pt}

\cvitem{21221学期}{文科高等数学 \quad 学生111人, 共64学时。
\newline 学生给出的量化评教平均分为: 99.0 (满分100)。}

\end{document}